\section{Uvod}

Informacioni sistem studija za animirane filmove predstavlja softversko rešenje namenjeno olakšavanju organizacije i praćenja svih faza produkcije animiranih filmova. Sistem omogućava vođenje evidencije o zaposlenima i njihovim ulogama, upravljanje filmskim projektima, budžetima, rokovima i zadacima, kao i praćenje materijala nastalih tokom produkcije. Takođe, obuhvata organizaciju rada produkcionih timova, pregled i odobravanje materijala, kao i upravljanje distribucijom završenih filmova. Cilj sistema je unapređenje efikasnosti poslovanja studija, bolja koordinacija zaposlenih i centralizovano čuvanje informacija o filmovima i njihovoj produkciji.

Projekat je realizovan u okviru kursa Informacioni sistemi na master studijama Matematičkog fakulteta Univerziteta u Beogradu.


\subsection{Akteri}
Sistem prepoznaje sledeće grupe korisnika koji interaktuju sa funkcionalnostima informacionog sistema na različite načine:


\begin{itemize}
	\item \textbf{Administrator}: Upravlja korisničkim nalozima, unosom i brisanjem zaposlenih, kao i njihovim dozvolama i pravima pristupa.
	\item \textbf{Producent}: Zadužen za kreiranje i ugovaranje projekata, upravljanje budžetima i rokovima, formiranje produkcionih timova i pregled izveštaja.
	\item \textbf{Projektni menadžer}: Organizuje rad na filmu, kreira i dodeljuje zadatke ostalim članovima tima, prati njihov napredak, pregleda i odobrava materijale i vraća ih na doradu po potrebi.
	\item \textbf{Član produkcionog tima}: Zajednički naziv za zaposlene koji izvršavaju zadatke na filmu: scenaristu, storyboardera, dizajnera, animatora, audio dizajnera i montažera. Dobija zadatke, predaje materijale i prati komentare i zahteve za doradu.
	\item \textbf{Animator}: Zadužen za kreiranje animacija na osnovu dodeljenih zadataka, predaju kreiranih materijala i ažuriranje njihovog statusa, kao i praćenje komentara i zahteva za doradu.
	\item \textbf{Scenarista}: Zadužen za kreiranje scenarija i praćenje njihovih verzija.
	\item \textbf{Storyboarder}: Pretvara scenario u niz vizuelnih kadrova, definiše njihov izgled, uglove kamere i tok radnje.
	\item \textbf{Dizajner}: Kreira i održava vizuelne materijale, poput dizajna likova i okruženja.
	\item \textbf{Audio dizajner}: Zadužen za dodavanje i obradu zvuka, muzike i zvučnih efekata.
	\item \textbf{Montažer}: Organizuje i povezuje kreirane materijale u finalnu verziju filma, integrišući zvuk i sliku.
	\item \textbf{Tehničko osoblje}: Zaduženo za održavanje i nabavljanje potrebne opreme.
	\item \textbf{Distributer}: Zadužen za plasiranje završenih filmova, ugovaranje njihovog prikazivanja, prodaje i distribucije, kao i za evidenciju distributivnih ugovora, licenci i prava prikazivanja.
\end{itemize}

\subsection{Korišćeni dijagrami i alati}
Za analizu i dizajn informacionog sistema korišćen je skup standardizovanih dijagrama i jezika za modelovanje, kao i odgovarajućih softverskih alata:

\begin{itemize}
	\item \textbf{UML (Unified Modeling Language)} dijagrami:
	\begin{itemize}
		\item \textbf{Slučajevi upotrebe}: Za prikaz funkcionalnosti sistema iz perspektive korisnika.
		\item \textbf{Dijagrami aktivnosti}: Za modelovanje toka radnji i poslovnih procedura.
		\item \textbf{Dijagrami sekvence}: Za prikaz interakcije između aktera i delova sistema tokom realizacije slučajeva upotrebe.
		\item \textbf{Dijagrami stanja}: Za prikaz životnog ciklusa važnih objekata u sistemu.
		\item \textbf{Dijagrami klasa}: Za prikaz modela baze podataka.
	\end{itemize}
	\item \textbf{BPMN (Business Process Model and Notation)}: Za detaljno modelovanje poslovnih procesa.
	\item \textbf{DFD (Data Flow Diagram)}: Za prikaz toka podataka kroz sistem.
	\item \textbf{LaTeX}: Korišćen kao alat za izradu i formatiranje tehničke dokumentacije.
\end{itemize}